\documentclass{article}
\usepackage[utf8]{inputenc}
\usepackage[T1]{fontenc}
\usepackage{lmodern}
\usepackage[a4paper, margin=1in]{geometry}
\usepackage[polish]{babel}

\usepackage{booktabs,tabularx,array}
\usepackage{graphicx}
\usepackage[table]{xcolor}
\usepackage[most]{tcolorbox}
\usepackage{xltabular}
\usepackage{enumitem}

% --- dodatkowe pakiety (diagramy, tabele hierarchiczne, odnośniki) ---
\usepackage{multirow}
\usepackage{float}
\usepackage{tikz}
\usetikzlibrary{positioning,arrows.meta,fit,backgrounds,calc,shapes.geometric,shapes.misc,babel}
\usepackage[hidelinks]{hyperref}

% --- kolory: poziomy hierarchii w tabelach konfiguracji i elementy diagramów ---
\definecolor{lvlA}{RGB}{176,196,222}   % poziom 1: sekcja pliku
\definecolor{lvlB}{RGB}{208,222,238}   % poziom 2: grupa
\definecolor{lvlC}{RGB}{232,240,248}   % poziom 3: typ metody
\definecolor{cPlik}{RGB}{222,235,250}
\definecolor{cSkrypt}{RGB}{255,232,204}
\definecolor{cZrodlo}{RGB}{220,242,220}
\definecolor{cWynik}{RGB}{245,225,245}
\definecolor{cModul}{RGB}{240,240,240}

\newcommand{\code}[1]{\texttt{#1}}

% --- bloki z poleceniami / fragmentami plików (tylko ASCII i polskie litery przez literate) ---
\newtcblisting{konsola}[1][]{
  listing only, colback=gray!5, colframe=gray!55, boxrule=0.4pt, arc=1mm,
  left=2mm, right=2mm, top=0.5mm, bottom=0.5mm,
  listing options={basicstyle=\ttfamily\footnotesize, columns=fullflexible,
    keepspaces=true, breaklines=true, showstringspaces=false,
    literate={ą}{{\k{a}}}1 {ć}{{\'c}}1 {ę}{{\k{e}}}1 {ł}{{\l{}}}1 {ń}{{\'n}}1
             {ó}{{\'o}}1 {ś}{{\'s}}1 {ź}{{\'z}}1 {ż}{{\.z}}1
             {Ą}{{\k{A}}}1 {Ć}{{\'C}}1 {Ę}{{\k{E}}}1 {Ł}{{\L{}}}1 {Ń}{{\'N}}1
             {Ó}{{\'O}}1 {Ś}{{\'S}}1 {Ź}{{\'Z}}1 {Ż}{{\.Z}}1},
  #1}

% --- style diagramów ---
\tikzset{
  >={Stealth[length=2.2mm]},
  plik/.style={draw=blue!50!black, fill=cPlik, rounded corners=1pt, font=\ttfamily\scriptsize,
               align=center, minimum height=8mm, inner sep=3pt},
  skrypt/.style={draw=orange!60!black, fill=cSkrypt, rounded corners=3mm, font=\scriptsize,
                 align=center, minimum height=8mm, inner sep=3pt},
  zrodlo/.style={draw=green!40!black, fill=cZrodlo, dashed, font=\scriptsize, align=center,
                 minimum height=8mm, inner sep=3pt},
  wynik/.style={draw=violet!60!black, fill=cWynik, font=\scriptsize, align=center,
                minimum height=8mm, inner sep=3pt, chamfered rectangle, chamfered rectangle xsep=2pt},
  modul/.style={draw=black!60, fill=cModul, font=\scriptsize, align=center,
                minimum height=7mm, inner sep=3pt},
  krok/.style={draw=black!70, fill=white, rounded corners=1pt, font=\scriptsize, align=center,
               minimum height=7mm, inner sep=3pt, text width=#1},
  krok/.default=4.2cm,
  petla/.style={draw=black!50, dashed, rounded corners=2mm, inner sep=3mm},
  strz/.style={->, thick, black!70},
  etyk/.style={font=\tiny\itshape, black!70, align=center},
  warstwa/.style={draw=black!25, fill=black!3, rounded corners=2mm, inner sep=2.5mm},
}

\large
\title{Projekt Zespołowy: Laboratorium 3}
\begin{document}
% babel-polish zamienia "c, "n, "o itd. na polskie litery; w tekście dokumentu cudzysłów
% występuje tylko w wartościach TOML (np. "classification"), więc skrót jest wyłączony
\shorthandoff{"}
\begin{titlepage}
	\begin{center}
    \line(1,0){300}\\
    [0.65cm]
	\huge{\bfseries Analiza porównawcza metod wczesnego zakończenia inferencji w modelach zespołowych}\\
	\line(1,0){300}\\
	\textsc{\Large Dokument modelowania}\\
	\textsc{\Large Projekt zespołowy: laboratorium 3}\\
    \textsc{\Large wersja 1.0}\\
	\textsc{\LARGE \today}\\
	[5.5cm]
	\end{center}
	\begin{flushright}
		\textsc{\Large Wojciech Baraniak}\\
		\textsc{\Large Maciej Winkler}\\
		[1cm]
        \textsc{\Large Promotor: mgr inż.\ Stanisław Kaźmierczak}\\
        \textsc{\Large Opiekun naukowy: prof.\ dr hab.\ inż.\ Jacek Mańdziuk}\\
		[1cm]
	\end{flushright}
\end{titlepage}
\newpage
\tableofcontents

\newpage

% =====================================================================================
\section*{Abstrakt}
\addcontentsline{toc}{section}{Abstrakt}

Las losowy złożony z setek drzew przy klasycznej predykcji odpytuje każde drzewo, choć
dla większości próbek wynik głosowania jest przesądzony już po kilkudziesięciu głosach.
Projekt bada metody \emph{wczesnego zakończenia inferencji} (ang.\ \emph{early stopping},
\emph{dynamic pruning}), które przerywają głosowanie, gdy wynik pełnego zespołu jest
wystarczająco pewny, i mierzy, ile pracy i czasu pozwalają zaoszczędzić oraz ile kosztuje
to w jakości i zgodności z pełnym lasem.

Niniejszy dokument opisuje stan implementacji po trzecim etapie projektu. Zaimplementowano
cztery metody z literatury (statystyczny \emph{instance-based pruning}, HYPER, GLEE z
systemu COMET oraz metodę opartą na marginesie lasu), metodę własną (HYPER z a priori
mieszanym z rozkładem jednostajnym, opcjonalnie z regułą GLEE), warianty regresyjne
oraz regułę wieloklasową ,,lider kontra druga klasa''. Całość badania jest sterowana jednym
plikiem konfiguracyjnym w formacie TOML: skrypt benchmarku uczy lasy na wygenerowanych
wcześniej zbiorach, kalibruje metody, mierzy predykcje w scenariuszach \emph{batch}
i \emph{online} i zapisuje surowe wyniki do pliku Parquet, z którego osobny moduł
analityczny liczy metryki i testy statystyczne.
Dokument przedstawia architekturę kodu, przepływ danych i plików, format konfiguracji,
procedury uruchomieniowe, parametry przykładowego pełnego badania (21 zbiorów
syntetycznych, 10 powtórzeń, 15 metod), listę prac, które pozostały do wykonania, oraz
budowę modelu: ustawienia lasu losowego i sposób głosowania drzew.

% =====================================================================================
\section*{Historia zmian}
\addcontentsline{toc}{section}{Historia zmian}

\begin{table}[H]
\centering
\small
\begin{tabularx}{\textwidth}{@{}llp{3.6cm}X@{}}
\toprule
\textbf{Wersja} & \textbf{Data} & \textbf{Autorzy} & \textbf{Opis zmian} \\
\midrule
1.0 & 18.11.2026 & W.~Baraniak, M.~Winkler &
Pierwsza wersja dokumentu: odniesienia do etapów 1 i 2, opis implementacji modułów,
diagramy przepływu danych, format konfiguracji, procedury uruchomieniowe, parametry
przykładowego badania, lista prac do wykonania, budowa modeli (parametry lasu,
soft voting a hard voting). \\
\bottomrule
\end{tabularx}
\end{table}

% =====================================================================================
\section{Odniesienia do dokumentów z etapów 1 i 2}

\subsection{Dokument wymagań (etap 1)}

Dokument wymagań określił cel projektu: porównanie metod wczesnego zakończenia
inferencji w zespołach drzew decyzyjnych pod względem kosztu predykcji (liczba
odpytanych drzew, czas) oraz jakości i zgodności z predykcją pełnego zespołu.
Tabela~\ref{tab:wymagania} zestawia główne wymagania z ich realizacją w kodzie.

\begin{table}[H]
\centering
\footnotesize
\caption{Realizacja wymagań z etapu 1.}
\label{tab:wymagania}
\begin{tabularx}{\textwidth}{@{}>{\raggedright\arraybackslash}p{4.4cm}X>{\centering\arraybackslash}p{1.7cm}@{}}
\toprule
\textbf{Wymaganie} & \textbf{Realizacja} & \textbf{Stan} \\
\midrule
Implementacja metod wczesnego stopu z literatury &
IB pruning, HYPER, GLEE, MARGIN w pakiecie \code{src/methods/} &
zrealizowane \\
Metody odniesienia &
\code{rf\_sklearn} (\code{rf.predict}, soft voting) i \code{full\_voting} (hard voting
wszystkich drzew); w regresji \code{full\_averaging} &
zrealizowane \\
Zbiory syntetyczne o kontrolowanych własnościach &
10 binarnych, 11 wieloklasowych, 8 regresyjnych (\code{generate\_datasets.py}) &
zrealizowane \\
Zbiory rzeczywiste &
SUSY, Statlog Shuttle, Gas Turbine pobierane z UCI (\code{data\_download\_real.py});
wstępna eksploracja w notatnikach \code{src/analysis/analysis\_*\_real.ipynb} &
pobrane, nie zbadane w benchmarku \\
Dwa scenariusze użycia &
\emph{batch} (cały zbiór testowy naraz) i \emph{online} (próbki pojedynczo) &
zrealizowane \\
Miary jakości, kosztu i zgodności &
accuracy, precision/recall/F1, MCC, \code{agree\_full}, \code{agree\_rf}, liczba drzew na
próbkę, czas, przyspieszenie (\code{src/analysis}, \code{analysis/binary.py}) &
zrealizowane \\
Powtarzalność eksperymentów &
wszystkie losowości sterowane ziarnami z konfiguracji; wyniki z metadanymi przebiegu &
zrealizowane \\
Testy statystyczne porównań &
test dwumianowy, McNemar, TOST, Friedman + Nemenyi, Wilcoxon (\code{analysis/binary.py}) &
tylko klasyfikacja binarna \\
\bottomrule
\end{tabularx}
\end{table}

\subsection{Analiza problemu i projekt rozwiązania (etap 2)}

Dokument projektu z etapu 2 zaproponował, poza czterema metodami z literatury,
rozszerzenia, które zostały zaimplementowane w tym etapie:

\begin{itemize}[nosep]
  \item \textbf{metoda własna} (\code{own}): HYPER z a priori
        $\lambda P_{\mathrm{OOB}} + (1-\lambda) P_{\mathrm{jedn}}$, opcjonalnie połączona
        z regułą GLEE (stop, gdy pozwala którakolwiek reguła);
  \item \textbf{warianty regresyjne} IB pruning i GLEE (\code{ib\_reg}, \code{glee\_reg})
        oraz regresyjna metoda własna z a priori wariancji wyjść drzew (\code{own\_reg});
        reguły stopu sprawdzono z bezpośrednią ewaluacją wzorów (2), (6) i (10)--(11)
        dokumentu projektu;
  \item \textbf{reguła pary} dla wielu klas (\code{multiclass = "pair"}): porównanie
        lidera z drugą klasą na tablicach binarnych;
  \item \textbf{kolejność drzew}: stała, jedna losowa permutacja lasu
        (\code{shuffled}) oraz osobna permutacja dla każdej próbki (\code{random});
  \item \textbf{kalibracja na zbiorze walidacyjnym} (\code{calib = "val"}) zamiast na
        głosach out-of-bag dla HYPER, metody własnej i MARGIN.
\end{itemize}

Elementy projektu, które nie zostały jeszcze zrealizowane, zebrano w
podrozdziale~\ref{sec:todo}.

% =====================================================================================
\section{Implementacja modułów}

\subsection{Działający kod}

\paragraph{Technologia.}
Kod jest napisany w Pythonie ($\ge 3.12$) na bibliotekach \code{numpy}, \code{scipy},
\code{pandas}, \code{pyarrow} i \code{scikit-learn}; analiza używa dodatkowo
\code{matplotlib} i \code{seaborn}. Las losowy to \code{RandomForestClassifier} lub
\code{RandomForestRegressor} ze scikit-learn. Wersje pakietów są przypięte w
\code{requirements.txt} i \code{requirements-dev.txt} (m.in.\ scikit-learn 1.9.1,
numpy 2.5.4, pandas 3.0.6). Funkcje i klasy mają adnotacje typów,
a repozytorium zawiera zestaw testów \code{pytest} (81 przypadków, wszystkie przechodzą).

\paragraph{Struktura repozytorium.}
Rysunek~\ref{fig:struktura} pokazuje strukturę katalogu \code{src}. Kod dzieli się na
cztery warstwy (konfiguracja, metody, benchmark, analiza), które komunikują się wyłącznie
przez jasno określone obiekty: konfigurację czyta tylko skrypt \code{testing.py} (pakiety
\code{methods} i \code{benchmarking} jej nie znają), a benchmark nie liczy żadnych metryk.

\begin{figure}[H]
\centering
\begin{tikzpicture}[
    every node/.style={anchor=west, font=\ttfamily\small, inner sep=2pt},
    opis/.style={font=\footnotesize\itshape, black!65},
    galaz/.style={draw=black!50}]
  \node at (0,0) {\textbf{src/}};
  \node at (0.6,-0.6) {config.py};       \node[opis] at (4,-0.6) {wczytanie i walidacja konfiguracji};
  \node at (0.6,-1.2) {testing.py};      \node[opis] at (4,-1.2) {skrypt benchmarku};
  \node at (0.6,-1.8) {methods/};        \node[opis] at (4,-1.8) {metody inferencji};
  \node at (0.6,-2.4) {benchmarking/};   \node[opis] at (4,-2.4) {trening lasu, pomiar, zapis wyników};
  \node at (0.6,-3.0) {data\_generate/}; \node[opis] at (4,-3.0) {generowanie i pobieranie zbiorów};
  \draw[galaz] (0.25,-0.25) -- (0.25,-3.0);
  \foreach \y in {-0.6,-1.2,-1.8,-2.4,-3.0} \draw[galaz] (0.25,\y) -- (0.6,\y);
\end{tikzpicture}
\caption{Struktura katalogu \code{src}.}
\label{fig:struktura}
\end{figure}

\paragraph{Przepływ danych -- widok ogólny.}
Rysunek~\ref{fig:przeglad} pokazuje, skąd pochodzą pliki wejściowe benchmarku. Kolorem
zielonym (linia przerywana) oznaczono źródła zewnętrzne, pomarańczowym skrypty, a
niebieskim pliki na dysku. Dalszy przepływ, od wczytania tych plików do zapisu wyników,
pokazuje rysunek~\ref{fig:benchmark}.

\begin{figure}[H]
\centering
\begin{tikzpicture}[node distance=7mm and 4mm]
  % źródła
  \node[zrodlo, text width=3.6cm] (uci) {Repozytorium UCI\\ \tiny archive.ics.uci.edu};
  \node[zrodlo, text width=3.6cm, right=of uci] (skl) {Generatory scikit-learn\\ \tiny make\_classification, make\_regression, make\_friedman1};
  \node[zrodlo, text width=3.6cm, right=of skl] (wzor) {Model urny / rekurencja\\ \tiny (bez danych, tylko $T$, $l$, $\pi$)};
  % skrypty danych
  \node[skrypt, text width=3.6cm, below=of uci] (dl) {data\_download\_real.py};
  \node[skrypt, text width=3.6cm, below=of skl] (gen) {generate\_datasets.py\\ \tiny ziarno: \code{seed} z config.toml};
  \node[skrypt, text width=3.6cm, below=of wzor] (gt) {generate\_table()};
  % pliki danych
  \node[plik, text width=3.6cm, below=of dl] (real) {data/real\_data/\\ *.parquet};
  \node[plik, text width=3.6cm, below=of gen] (gend) {data/generated/\\ *.parquet};
  \node[plik, text width=3.6cm, below=of gt] (tab) {data/tablice/\\ T\{T\}\_L\{l\}\_pi\{pi\}.npz};

  \draw[strz] (uci) -- (dl);
  \draw[strz] (skl) -- (gen);
  \draw[strz] (wzor) -- (gt);
  \draw[strz] (dl) -- (real);
  \draw[strz] (gen) -- (gend);
  \draw[strz] (gt) -- (tab);
\end{tikzpicture}
\caption{Pochodzenie plików wejściowych benchmarku: zbiorów danych i tablic progów IB pruning.}
\label{fig:przeglad}
\end{figure}

\paragraph{Zbiory danych.}
Wszystkie zbiory zapisywane są w jednolitym formacie Parquet: kolumny cech
\code{f0}, \code{f1}, \dots, kolumna celu \code{target} oraz ewentualne kolumny
dodatkowe, które nie są cechami (w \code{gas\_turbine}: \code{year} i drugi cel
\code{CO}, pomijane dzięki \code{benchmark.drop\_columns}).

\paragraph{Przebieg benchmarku.}
Skrypt \code{src/testing.py} realizuje pętlę po zbiorach danych i ziarnach (rysunek~\ref{fig:benchmark}).
Dla każdej pary (zbiór, ziarno) dzieli dane, uczy jeden las, wspólny dla wszystkich
metod, tworzy \emph{nowe} instancje metod i kalibruje je na tym lesie, mierzy predykcje
w wybranych scenariuszach i od razu dopisuje wyniki do pliku. Zbiór testowy jest wspólny
dla wszystkich metod, więc porównania są parami (ang.\ \emph{paired}).

\begin{figure}[H]
\centering
\begin{tikzpicture}[node distance=4.5mm]
  \node[plik, text width=4.2cm] (cfg) {config.toml $\to$ config.py\\ \tiny stałe TASK, DATASETS, SEEDS, MODES, METHODS, \dots};
  \node[krok, below=of cfg] (bm) {\code{build\_methods(METHODS)}\\ \tiny walidacja konfiguracji, lista metod};
  \node[krok, below=of bm] (fd) {\code{find\_datasets(DATASETS)}\\ \tiny folder $\to$ wszystkie pliki *.parquet};
  \node[krok, below=11mm of fd] (ld) {\code{load\_dataset(path, DROP\_COLUMNS)} $\to X, y$};
  \node[krok, below=11mm of ld] (sp) {\code{train\_test\_split}: test (\code{test\_size})\\ \tiny opcjonalnie drugi podział: walidacja (\code{val\_size})};
  \node[krok, below=of sp] (tf) {\code{train\_forest}\\ \tiny $n$\_trees, max\_depth (z nadpisaniami per zbiór), ziarno};
  \node[krok, below=of tf] (kal) {\code{ForestMethods.add\_method(m)}\\ $\to$ \code{m.fit(rf, X\_train, y\_train, X\_val, y\_val)}};
  \node[krok, below=of kal] (run) {\code{Benchmark.run(MODES)}\\ \tiny batch: cały \code{X\_test}; online: pierwsze \code{n\_online} próbek;\\ czas = mediana z \code{n\_repeats}};
  \node[krok, below=of run] (rel) {\code{release()} -- zwolnienie lasu i danych};
  \node[krok, below=of rel] (add) {\code{BenchmarkSuite.add} $\to$ \code{store\_runs}};
  \node[plik, below=8mm of add, text width=4.2cm] (res) {data/results/\{out\_file\}.parquet};

  % źródła dla fit
  \node[plik, right=12mm of kal, text width=3.6cm] (npz) {data/tablice/*.npz\\ \tiny (IB pruning: dobór wg $T$, $l$, $\pi$)};
  \node[modul, right=12mm of tf, text width=3.6cm] (oob) {głosy out-of-bag lasu\\ \tiny HYPER, OWN, MARGIN (\code{calib="oob"})};
  \node[modul, right=12mm of sp, text width=3.6cm] (val) {\code{X\_val}, \code{y\_val}\\ \tiny HYPER, OWN, MARGIN (\code{calib="val"})};
  \node[modul, left=12mm of kal, text width=3.2cm] (glee) {tylko $T$ i $\pi$\\ \tiny GLEE: tablica liczona w \code{fit}};
  \node[plik, left=12mm of ld, text width=3.2cm] (data) {data/generated/\dots\\ data/real\_data/\dots};
  \node[modul, left=12mm of run, text width=3.2cm] (tryb) {\code{predict\_batch(X)}\\ \code{predict\_one(x)}\\ \tiny $\to$ (y\_pred, n\_used)};

  \draw[strz] (cfg) -- (bm);
  \draw[strz] (bm) -- (fd);
  \draw[strz] (fd) -- (ld);
  \draw[strz] (ld) -- (sp);
  \draw[strz] (sp) -- (tf);
  \draw[strz] (tf) -- (kal);
  \draw[strz] (kal) -- (run);
  \draw[strz] (run) -- (rel);
  \draw[strz] (rel) -- (add);
  \draw[strz] (add) -- node[etyk, right]{zapis od razu\\ po przebiegu} (res);
  \draw[strz, densely dashed] (npz.west) -- ([yshift=-2mm]kal.east);
  \draw[strz, densely dashed] (oob.south) -- ++(0,-1mm) -| ([xshift=10mm]kal.north);
  \draw[strz, densely dashed] (val.west) -- ++(-4mm,0) |- ([yshift=2.5mm]kal.east);
  \draw[strz, densely dashed] (glee) -- (kal);
  \draw[strz, densely dashed] (data) -- (ld);
  \draw[strz, densely dashed] (tryb) -- (run);

  \begin{scope}[on background layer]
    \node[petla, fit=(sp)(add), label={[font=\scriptsize\bfseries, anchor=south west, xshift=1mm]north west:dla każdego ziarna z \code{seeds}}] (p2) {};
    \node[petla, fit=(ld)(p2), inner sep=4mm, label={[font=\scriptsize\bfseries, anchor=south west, xshift=1mm]north west:dla każdego pliku z \code{datasets}}] {};
  \end{scope}
\end{tikzpicture}
\caption{Przebieg skryptu \code{src/testing.py}. Linie przerywane: dane, z których
korzystają metody przy kalibracji (\code{fit}) i pomiarze.}
\label{fig:benchmark}
\end{figure}

Zbiór walidacyjny jest odkładany z danych uczących \emph{przed} treningiem lasu, więc
las go nie widzi; \code{val\_size} jest udziałem w \emph{całym} zbiorze, dlatego przy
\code{test\_size = 0.3} i \code{val\_size = 0.1} las uczy się na 60\% danych.

\paragraph{Metody inferencji.}
Wszystkie metody implementują wspólny kontrakt \code{InferenceMethod}:
\code{fit(rf, X\_train, y\_train, X\_val, y\_val)}, \code{predict\_batch(X)} oraz
\code{predict\_one(x)}; obie metody predykcji zwracają parę (predykcje, liczba
odpytanych drzew na próbkę). Klasy pośrednie zawierają wspólne pętle głosowania
(\code{FixedVoting}, \code{SequentialPruning}, a w regresji \code{SequentialRegression}),
więc nowa metoda podaje jedynie regułę stopu albo wybór drzew (rysunek~\ref{fig:klasy}).
Drzewa są ewaluowane przez lekką klasę \code{FastTree} (zejście po drzewie i odczyt klasy
liścia bez walidacji wejścia scikit-learn; w regresji \code{FastRegTree}), co dla
pojedynczej próbki skraca wywołanie drzewa z ok.\ 68\,$\mu$s do kilku mikrosekund; dzięki
temu pomiar czasu mierzy metody, a nie narzut biblioteki. Metoda odniesienia
\code{rf\_sklearn} celowo używa oryginalnego \code{rf.predict}.

\begin{figure}[H]
\centering
\begin{tikzpicture}[
    kl/.style={modul, font=\ttfamily\scriptsize, minimum height=6mm},
    abs/.style={kl, fill=black!12, font=\ttfamily\itshape\scriptsize},
    dz/.style={-{Triangle[open, length=2.2mm]}, black!70}]
  % poziom 0
  \node[abs] (base) at (0,0) {InferenceMethod};
  % poziom 1
  \node[kl]  (skl)   at (-6.0,-1.3) {SklearnRF};
  \node[abs] (fixed) at (-3.4,-1.3) {FixedVoting};
  \node[abs] (seq)   at ( 0.3,-1.3) {SequentialPruning};
  \node[abs] (reg)   at ( 4.6,-1.3) {RegressionMethod};
  % poziom 2
  \node[kl] (full) at (-4.6,-2.6) {FullVoting};
  \node[kl] (marg) at (-2.3,-2.6) {MarginEarlyStop};
  \node[kl] (fa)   at ( 2.6,-2.6) {FullAveraging};
  \node[kl] (sklr) at ( 6.3,-2.6) {SklearnRFReg};
  % poziom 3
  \node[kl]  (ib)   at (-1.9,-3.9) {IBPruning};
  \node[kl]  (hyp)  at ( 0.3,-3.9) {Hyper};
  \node[kl]  (glee) at ( 2.1,-3.9) {Glee};
  \node[abs] (sr)   at ( 5.0,-3.9) {SequentialRegression};
  % poziom 4
  \node[kl] (own) at (-0.6,-5.2) {Own};
  \node[kl] (ibr) at ( 2.2,-5.2) {IBRegression};
  \node[kl] (gr)  at ( 4.5,-5.2) {GleeRegression};
  \node[kl] (or)  at ( 6.8,-5.2) {OwnRegression};
  % dziedziczenie (strzałka do klasy bazowej)
  \foreach \dziecko/\rodzic in {skl/base, fixed/base, seq/base, reg/base,
                     full/fixed, marg/fixed, fa/reg, sklr/reg,
                     ib/seq, hyp/seq, glee/seq, sr/reg,
                     own/hyp, ibr/sr, gr/sr, or/sr}
    \draw[dz] (\dziecko.north) -- (\rodzic.south);
  \draw[dz, densely dashed] (own.north) -- (ib.south);
\end{tikzpicture}
\caption{Hierarchia klas metod (\code{src/methods/}). Kursywą oznaczono klasy abstrakcyjne:
\code{FixedVoting} -- stała lista drzew dla każdej próbki, \code{SequentialPruning}
i \code{SequentialRegression} -- głosowanie sekwencyjne z regułą stopu. Linia przerywana:
\code{Own} łączy a priori HYPER z jednostajnym a priori IB pruning (waga $\lambda$).}
\label{fig:klasy}
\end{figure}

\begin{table}[H]
\centering
\footnotesize
\caption{Zaimplementowane metody i źródła danych, z których korzystają przy kalibracji.}
\label{tab:metody}
\begin{tabularx}{\textwidth}{@{}>{\ttfamily}lX>{\raggedright\arraybackslash}p{4.3cm}@{}}
\toprule
\normalfont\textbf{Typ} & \textbf{Reguła} & \textbf{Kalibracja w \code{fit}} \\
\midrule
rf\_sklearn & \code{rf.predict} ze scikit-learn (soft voting) & brak \\
full\_voting & hard voting wszystkich $T$ drzew (odniesienie dla \code{agree\_full}) & brak \\
ib\_pruning & stop, gdy $P(\text{lider wygra po } T \text{ głosach}) \ge \pi$ przy a priori jednostajnym (model urny) & tablica progów z pliku \code{.npz} \\
hyper & jak IB, a priori rozkładu głosów szacowane z danych & głosy OOB lub walidacyjne \\
glee & dolna granica przedziału ufności udziału lidera $> 0{,}5$ (poprawka na populację skończoną) & tablica z $T$ i $\pi$ \\
margin & pierwsze $N^* = \lceil T(\mu + k\sigma)\rceil$ drzew rankingu MeanD-M & margines na danych OOB lub walidacyjnych \\
own & HYPER z a priori $\lambda P_{\mathrm{OOB}} + (1-\lambda)P_{\mathrm{jedn}}$, opcjonalnie lub GLEE & głosy OOB lub walidacyjne \\
\midrule
full\_averaging & średnia wszystkich drzew (regresja) & brak \\
ib\_reg & przedział $t$-Studenta dla średniej pełnego lasu w $\pm\delta$ & tylko $T$, $\pi$, $\delta$ \\
glee\_reg & przedział normalny z poprawką na populację skończoną & tylko $T$, $\pi$, $\delta$ \\
own\_reg & jak \code{ib\_reg}, z a priori wariancji wyjść drzew ($\nu_0$ pseudoobserwacji) & predykcje OOB \\
\bottomrule
\end{tabularx}
\end{table}

\paragraph{Format pliku wyników.}
Benchmark nie liczy metryk: zapisuje do jednego pliku Parquet jeden wiersz na
przebieg $\times$ scenariusz $\times$ metodę, z ustawieniami przebiegu, czasem predykcji
i surowymi tablicami (\code{y\_true}, \code{y\_pred}, \code{n\_used}). Metryki liczy
później moduł \code{src/analysis}, dlatego poprawka lub nowa metryka obejmuje także
wcześniej policzone przebiegi.

\paragraph{Analiza wyników.}
Druga, niezależna od benchmarku część potoku to moduł \code{src/analysis}:
\code{compute\_metrics} zamienia surowe tablice na metryki: accuracy, makro
precision/recall/F1, zgodność z \code{rf\_sklearn} i \code{full\_voting}
(\code{agree\_rf}, \code{agree\_full}), czas na próbkę, przyspieszenie oraz liczbę
odpytanych drzew (\code{trees\_mean}, \code{trees\_median}, \code{trees\_max},
\code{trees\_frac}); w regresji MAE, RMSE, $R^2$ oraz średnią i maksymalną różnicę od
pełnego lasu. Skrypt \code{analysis/binary.py} dokłada miary dla klasy mniejszościowej,
MCC, zrównoważoną dokładność, zgodność w podziale na klasy, koszt na klasie rzadkiej
i częstej oraz testy statystyczne.

\paragraph{Plik konfiguracyjny i parametry metod.}
Badanie jest w całości opisane plikiem TOML (domyślnie \code{config.toml} w katalogu głównym,
inny plik można podać jako argument skryptu). Sekcja \code{[benchmark]} wybiera zadanie
i zbiory, podział danych (\code{seeds}, \code{test\_size}, \code{val\_size}) oraz
scenariusze pomiaru (\code{modes}, \code{n\_online}, \code{n\_repeats}); sekcja
\code{[forest]} podaje parametry lasu z możliwością nadpisania dla wybranych zbiorów
(\code{[forest.per\_dataset.<zbiór>]}); każdy blok \code{[[methods]]} opisuje jedną badaną
metodę.

Parametry metod podaje się w blokach \code{[[methods]]}; w nawiasach podano wartości
domyślne. Wszystkie metody z wczesnym stopem poza MARGIN przyjmują kolejność drzew
\code{order} (\code{fixed}, \code{shuffled} albo \code{random}) i ziarno permutacji
\code{seed}, a metody odniesienia \code{full\_voting} i \code{full\_averaging} nie mają
parametrów. Podstawowym parametrem metod probabilistycznych jest poziom ufności $\pi$
(\code{pi}, 0,99): IB pruning wybiera na jego podstawie plik tablicy progów
\code{T\{T\}\_L\{l\}\_pi\{100$\pi$\}.npz}, a GLEE przyjmuje poziom istotności
$\alpha = 1 - \pi$. Parametr \code{multiclass} (\code{full} dla 2--3 klas albo \code{pair}
dla dowolnej liczby klas) określa, czy reguła stopu działa na pełnym rozkładzie głosów,
czy tylko na parze lider -- druga klasa. HYPER szacuje a priori rozkładu głosów z
\code{n\_prior} (20\,000) próbek, wygładza je oknem \code{smooth} (5) w trybie
\code{smooth\_mode} (\code{mean} albo \code{mass}) i kalibruje się na głosach out-of-bag
albo na zbiorze walidacyjnym (\code{calib}: \code{oob} albo \code{val}). GLEE nie
zatrzymuje głosowania przed zebraniem \code{min\_votes} głosów (15, 30 lub 45 zależnie od
$\alpha$), stosuje poprawkę na populację skończoną (\code{fpc}) i przedział jednostronny
(\code{one\_sided}), a kolejność drzew domyślnie losuje osobno dla każdej próbki
(\code{order = "random"}, \code{seed = 0}). MARGIN dla każdej próbki odpytuje tę samą
liczbę pierwszych drzew rankingu (\code{ranking}: \code{mean\_d\_m}, \code{accuracy} albo
\code{none}):
\[
  N^* = \bigl\lceil T\,(\mu + k\sigma) \bigr\rceil ,
\]
gdzie $\mu$ i $\sigma$ to średnia i odchylenie standardowe marginesu poprawnie
sklasyfikowanych próbek kalibracyjnych (\code{n\_calib}, 20\,000), a $k$ (\code{k}, 0)
przesuwa liczbę drzew. Metoda własna przyjmuje parametry HYPER, a jej a priori jest
mieszaniną
\[
  P = \lambda\, P_{\mathrm{OOB}} + (1 - \lambda)\, P_{\mathrm{jedn}} ,
\]
gdzie $\lambda$ (\code{lam}, 0,5) równe 1 daje HYPER, a równe 0 IB pruning; przy
\code{glee = true} głosowanie kończy się, gdy pozwala reguła HYPER lub GLEE
(\code{pi\_glee}, domyślnie równe \code{pi}, oraz \code{min\_votes} jak w GLEE).

W regresji \code{ib\_reg} kończy głosowanie, gdy przedział ufności na poziomie $\pi$ dla
średniej pełnego lasu mieści się w $\pm\delta$. Tolerancję podaje się wprost (\code{delta})
albo względnie (\code{delta\_std}):
\[
  \delta = \mathtt{delta\_std} \cdot \sigma_{y_{\mathrm{train}}} ,
\]
a \code{t\_min} (2) to minimalna liczba drzew przed stopem. \code{glee\_reg} ma te same
parametry, \code{t\_min} równe 15, 30 lub 45 i poprawkę \code{fpc}. \code{own\_reg}
dodaje a priori wariancji wyjść drzew o sile $\nu_0$ pseudoobserwacji (\code{nu0}, 5;
$\nu_0 = 0$ daje \code{ib\_reg}), szacowane z \code{n\_prior} (20\,000) predykcji
out-of-bag, i dopuszcza stop już po pierwszym drzewie (\code{t\_min = 1}).

\paragraph{Uruchamianie kodu.}
Środowisko przygotowuje jedno polecenie \code{setup\_windows.bat} uruchomione z katalogu
głównego. Skrypt sprawdza, czy dostępny jest Python $\ge 3.12$, tworzy środowisko
\code{.venv} (lub używa istniejącego) i instaluje w nim przypięte pakiety z
\code{requirements-dev.txt}. Następnie generuje zbiory syntetyczne do
\code{data/generated/} i pobiera zbiory rzeczywiste z UCI do \code{data/real\_data/};
pobieranie jest pomijane, gdy pliki już istnieją, chyba że podano flagę
\code{--force-download}. W innych systemach te same kroki wykonuje się ręcznie
(\code{python -m venv}, \code{pip install} i oba skrypty z \code{src/data\_generate/}).
Badanie uruchamia skrypt \code{src/testing.py} z plikiem konfiguracyjnym jako argumentem,
metryki i podsumowanie liczy \code{analysis.summary}, a analizę porównawczą
\code{analysis/binary.py}. Nowe tablice progów IB pruning (dla innych $T$, $l$, $\pi$)
tworzy funkcja \code{methods.generate\_table}.

\paragraph{Przykładowe badanie: pełny przebieg klasyfikacji.}
Jako przykładowe badanie wykonano pełny przebieg klasyfikacji opisany plikiem
\code{configs/classification\_all.toml} (tabela~\ref{tab:badanie}); przebieg
uzupełniający (\code{configs/classification\_quick.toml}, 3 zbiory, ziarna 1--5) dodał
metodę \code{GLEE\_pi98\_rand}.

\begin{table}[H]
\centering
\footnotesize
\caption{Parametry przykładowego badania (\code{configs/classification\_all.toml}).}
\label{tab:badanie}
\begin{tabularx}{\textwidth}{@{}>{\raggedright\arraybackslash}p{4.2cm}X@{}}
\toprule
\textbf{Parametr} & \textbf{Wartość} \\
\midrule
Zadanie & klasyfikacja (\code{RandomForestClassifier}) \\
Zbiory & 21 syntetycznych: 10 binarnych i 11 wieloklasowych (3--20 klas); zbiory testowe od 300 do 150\,000 próbek \\
Las & $T = 501$ drzew, \code{max\_depth = 5}; dla \code{binary\_multi\_features} i \code{multi\_multi\_features} bez limitu głębokości \\
Podział danych & 60\% trening, 10\% walidacja (\code{val\_size = 0.1}), 30\% test \\
Powtórzenia & 10 ziaren (1--10): podział danych i las zależą od ziarna; razem 210 przebiegów \\
Scenariusze & \emph{batch} (cały zbiór testowy) i \emph{online} (pierwsze 500 próbek pojedynczo) \\
Pomiar czasu & mediana z 4 powtórzeń, jeden wątek \\
Odniesienia & \code{rf\_sklearn} (soft voting) i \code{full\_voting} (hard voting) \\
Poziom ufności & $\pi = 0{,}98$ dla wszystkich metod probabilistycznych; MARGIN z $k = 0$ \\
\bottomrule
\end{tabularx}
\end{table}

\subsection{Procedury automatyzacji}

\begin{itemize}[leftmargin=*]
  \item \textbf{Badanie sterowane konfiguracją.} Zmiana zbiorów, parametrów lasu, metod,
        ich wariantów i scenariuszy nie wymaga zmian w kodzie. Rejestr
        \code{METHOD\_TYPES} zamienia każdy blok \code{[[methods]]} na wywołanie
        konstruktora i waliduje parametry (oraz unikalność nazw metod) przed rozpoczęciem
        obliczeń, więc błąd w konfiguracji ujawnia się od razu, a nie po godzinach obliczeń.
  \item \textbf{Automatyczny dobór zasobów metod.} IB pruning wybiera plik tablicy progów
        pasujący do $T$ i $l$ danego lasu; metody, które nie pasują do lasu (np.\ wariant
        \code{full} przy 4+ klasach, \code{calib="val"} bez zbioru walidacyjnego), są
        pomijane z jednorazowym komunikatem, a reszta badania trwa dalej.
  \item \textbf{Zapis przyrostowy i odporność na przerwanie.} Każdy przebieg (zbiór,
        ziarno) jest dopisywany do pliku wyników od razu po zakończeniu, przez plik
        tymczasowy i atomową podmianę. Przerwany benchmark zachowuje policzone przebiegi,
        a nieudany zapis (np.\ plik otwarty w innym programie) nie przerywa obliczeń --
        przebieg jest zapisywany przy następnej próbie.
        Wiersz z tymi samymi ustawieniami, scenariuszem i metodą zastępuje zapisany, a
        przebiegi z innymi ustawieniami dopisują się obok -- jeden plik może więc
        gromadzić kolejne eksperymenty (np.\ różne \code{max\_depth}).
  \item \textbf{Testy automatyczne.} \code{python -m pytest} uruchamia 81 testów:
        poprawność reguł stopu (porównanie z bezpośrednią ewaluacją wzorów), zgodność
        trybów batch i online, przypadki skrajne (np.\ $\lambda = 0$ i $\lambda = 1$ w
        metodzie własnej), wczytywanie konfiguracji, rejestr metod, zapis i odczyt pliku
        wyników oraz metryki.
\end{itemize}

\subsection{Co jeszcze nie zostało wykonane}
\label{sec:todo}

\begin{table}[H]
\centering
\footnotesize
\begin{tabularx}{\textwidth}{@{}>{\raggedright\arraybackslash}p{4.6cm}X@{}}
\toprule
\textbf{Zadanie} & \textbf{Stan i uwagi} \\
\midrule
\multicolumn{2}{@{}l}{\textit{Eksperymenty}} \\
Ponowny pełny przebieg klasyfikacji &
wyniki sprzed 16.11.2026 są nieaktualne dla HYPER, OWN, GLEE i MARGIN (poprawki zgodności
z artykułami źródłowymi); aktualny jest tylko przebieg uzupełniający (3 zbiory, ziarna
1--5). Pełny przebieg trzeba powtórzyć, najlepiej do nowego pliku \\
Przebieg regresji &
metody i 8 zbiorów są gotowe, bloki są w \code{config.toml}; brak pliku wyników \\
Zbiory rzeczywiste &
SUSY, Shuttle i Gas Turbine są pobrane, ale nie były badane w benchmarku (SUSY: 5 mln
wierszy, wymaga oszacowania czasu na jednym ziarnie) \\
Dodatkowe zbiory &
surowe pliki MetroPT, WISDM i czujników gazu (\code{data/real\_data/extra\_time\_data})
nie są jeszcze przekształcone do formatu projektu \\
Wpływ parametrów &
badania dla różnych $\pi$ (wymaga nowych tablic IB), różnych $T$ i \code{max\_depth}
oraz różnych $k$ w MARGIN i $\lambda$ w metodzie własnej \\
\midrule
\multicolumn{2}{@{}l}{\textit{Analiza}} \\
Analiza wieloklasowa i regresyjna &
\code{analysis/} zawiera dotąd tylko skrypt dla klasyfikacji binarnej \\
Pełna analiza binarna &
\code{analysis/binary.py} był uruchamiany tylko na wycinku wyników sprzed poprawek
(ziarna 1--3); po ponownym przebiegu trzeba go uruchomić na pełnym pliku \\
\midrule
\multicolumn{2}{@{}l}{\textit{Elementy dokumentu projektu (etap 2)}} \\
Kalibracja do wspólnego poziomu niezgodności &
dobór $\pi$ / $k$ tak, by metody miały ten sam odsetek niezgodności na zbiorze
walidacyjnym -- nie zaimplementowano \\
Strojenie $\nu_0$ w \code{own\_reg} i $k$ w MARGIN &
obecnie wartości stałe z konfiguracji \\
Szybki ewaluator drzew (Numba) &
nie zaimplementowano; obecnie \code{FastTree} na \code{tree\_.apply} \\
Inne zespoły niż las losowy &
bagging, IVoting -- poza obecnym zakresem kodu \\
\bottomrule
\end{tabularx}
\end{table}

% =====================================================================================
\section{Budowa modeli}
\label{sec:modele}


Uczonym modelem jest las losowy ze scikit-learn (\code{RandomForestClassifier},
w regresji \code{RandomForestRegressor}). Dla każdego zbioru i ziarna powstaje jeden las,
wspólny dla wszystkich badanych metod, więc różnice między metodami wynikają wyłącznie
ze sposobu głosowania drzew.

\begin{itemize}[leftmargin=*]
  \item \textbf{Liczba drzew: $T = 501$} (\code{[forest] n\_trees}). Las tej wielkości
        uczy się i ewaluuje w rozsądnym czasie, a jednocześnie jest na tyle duży, że
        wczesny stop daje mierzalne oszczędności, a wyniki metod są porównywalne między sobą.
  \item \textbf{Głębokość drzew: \code{max\_depth = 5}} (\code{[forest] max\_depth}).
        Wyjątkiem są zbiory z wieloma cechami informatywnymi (\code{binary\_multi\_features},
        \code{multi\_multi\_features}: 100 cech, z czego 50 informatywnych), dla których
        \code{max\_depth} ma wartość domyślną scikit-learn, czyli brak limitu głębokości
        (\code{max\_depth = "none"} w \code{[forest.per\_dataset.<zbiór>]}). Przy
        głębokości 5 drzewa są na tych zbiorach zbyt słabe i las traci dokładność.
\end{itemize}

Drzewa są zawsze uczone na próbach bootstrapowych (\code{bootstrap=True}); jest to
warunek kalibracji na głosach out-of-bag. Pozostałe parametry lasu mają wartości
domyślne scikit-learn.

Las losowy scikit-learn w metodzie \code{predict} stosuje \emph{soft voting}: uśrednia
prawdopodobieństwa klas zwracane przez drzewa (rozkłady klas w liściach) i wybiera klasę
o największym średnim prawdopodobieństwie. Badane metody bazują natomiast na
\emph{twardych} odpowiedziach drzew (\emph{hard voting}): każde drzewo oddaje jeden głos
na klasę przeważającą w swoim liściu, a wynik zależy od liczby głosów. Tylko przy takim
głosowaniu można przerwać odpytywanie drzew, gdy przewaga głosów lidera przesądza wynik
pełnego zespołu.

Dlatego punktem odniesienia dla metod wczesnego stopu jest \code{full\_voting} --
hard voting wszystkich 501 drzew (miara \code{agree\_full}), a oryginalny \code{rf.predict}
(\code{rf\_sklearn}) jest dodatkowym odniesieniem, które pokazuje różnicę względem
domyślnej predykcji scikit-learn (miara \code{agree\_rf}).


\end{document}
